\chapter{CÀI ĐẶT VÀ KIỂM THỬ}

\section{Công cụ kiểm thử}

\subsection{Google Lighthouse}
Google Lighthouse là một công cụ kiểm thử và phân tích website 
mã nguồn mở do Google phát triển, được sử dụng rộng rãi trong 
việc đánh giá chất lượng và tối ưu hóa trang web \cite{lighthouse_docs}. 
Công cụ này có thể chạy trực tiếp trên Chrome DevTools, cài đặt 
dưới dạng tiện ích mở rộng hoặc sử dụng qua dòng lệnh. Khi thực 
hiện kiểm tra, Lighthouse tự động phân tích nhiều khía cạnh quan 
trọng của website bao gồm hiệu năng tải trang (Performance), khả 
năng truy cập (Accessibility), mức độ tuân thủ các chuẩn SEO, 
các thực hành tốt về bảo mật (Best Practices) và khả năng hỗ trợ 
Progressive Web App (PWA). Sau quá trình đánh giá, công cụ cung 
cấp báo cáo chi tiết kèm các chỉ số Core Web Vitals (FCP, LCP, 
CLS, TBT, Speed Index) và gợi ý cải thiện cụ thể.

\subsection{Postman}
Postman là công cụ kiểm thử API phổ biến, cho phép gửi các yêu 
cầu HTTP đến server và kiểm tra phản hồi một cách trực quan 
\cite{postman_docs}. Trong đề tài, Postman được sử dụng để kiểm 
thử các endpoint của FastAPI backend bao gồm xác thực JWT, upload 
ảnh cho AI Virtual Try-On, các API quản lý sản phẩm, đơn hàng, 
chatbot và thanh toán Stripe. Postman hỗ trợ kiểm thử với 
multipart/form-data, thiết lập header xác thực và lưu trữ 
collection test case có tổ chức.

\subsection{Swagger UI}
Swagger UI là giao diện tài liệu và kiểm thử API được FastAPI 
tự động sinh ra tại đường dẫn \texttt{/docs} \cite{fastapi_docs}. 
Công cụ này cho phép kiểm thử trực tiếp tất cả endpoint của hệ 
thống ngay trên trình duyệt mà không cần cài đặt thêm phần mềm, 
đồng thời hiển thị đầy đủ schema request và response của từng API.

\subsection{Responsive Viewer}
Responsive Viewer là tiện ích mở rộng (extension) của trình duyệt 
Chrome, cho phép kiểm tra giao diện website trên nhiều kích thước 
màn hình cùng lúc \cite{responsive_viewer}. Trong đề tài, công cụ 
này được sử dụng để kiểm thử tính tương thích giao diện của hệ 
thống trên các thiết bị di động, máy tính bảng và máy tính để bàn 
mà không cần thiết bị vật lý.

% ─────────────────────────────────────────
\section{Kiểm thử}

\subsection{Kiểm thử API}

Các endpoint của hệ thống được kiểm thử thông qua Postman và 
Swagger UI. Bảng~\ref{tab:api_test} trình bày một số test case 
tiêu biểu.

\begin{table}[H]
\centering
\caption{Một số test case kiểm thử API}
\label{tab:api_test}
\begin{tabular}{|c|p{3cm}|p{3cm}|p{3cm}|c|}
\hline
\textbf{STT} & \textbf{Endpoint} & \textbf{Input} & 
\textbf{Kết quả mong đợi} & \textbf{Kết quả} \\
\hline
1 & POST /users & Firebase UID hợp lệ & 
Upsert user thành công, trả về thông tin user & Đạt \\
\hline
2 & GET /garments/ & -- & 
Trả về danh sách garment & Đạt \\
\hline
3 & POST /tryon & Ảnh người + ảnh quần áo & 
Trả về URL ảnh kết quả từ Cloudinary & Đạt \\
\hline
4 & POST /photo-tryon/fit-assessment & 
body\_profile\_id, garment\_id, size\_label & 
Trả về đánh giá chest\_fit, waist\_fit, hip\_fit, 
size\_suggestion & Đạt \\
\hline
5 & POST /chatbot/chat & Câu hỏi tư vấn thời trang & 
Trả về phản hồi từ Ollama qwen2.5 & Đạt \\
\hline
6 & POST /payment/stripe & Thông tin đơn hàng & 
Tạo payment intent thành công & Đạt \\
\hline
7 & POST /tryon & Thiếu file ảnh & 
Trả về lỗi 422 Unprocessable Entity & Đạt \\
\hline
\end{tabular}
\end{table}

% ─────────────────────────────────────────
\subsection{Kiểm thử hiệu năng}

Kiểm thử hiệu năng được thực hiện bằng Google Lighthouse trên 
hai trang quan trọng của hệ thống là trang chủ và trang danh 
sách sản phẩm.

\begin{figure}[H]
    \centering
    \includegraphics[width=\textwidth]{graphics/test/lighthouse_home.png}
    \caption{Kiểm thử hiệu năng trang chủ}
    \label{fig:lighthouse_home}
\end{figure}

\begin{figure}[H]
    \centering
    \includegraphics[width=\textwidth]{graphics/test/lighthouse_product.png}
    \caption{Kiểm thử hiệu năng trang danh sách sản phẩm}
    \label{fig:lighthouse_product}
\end{figure}

Các chỉ số đo lường trong kết quả Lighthouse:

\begin{itemize}
    \item \textbf{First Contentful Paint (FCP)}: Thời gian từ khi 
    người dùng mở trang cho đến khi trình duyệt hiển thị nội dung 
    đầu tiên (văn bản, ảnh, canvas...).
    
    \item \textbf{Largest Contentful Paint (LCP)}: Thời gian để 
    hiển thị phần tử nội dung lớn nhất trong vùng nhìn thấy, 
    thường là ảnh sản phẩm hoặc banner.
    
    \item \textbf{Total Blocking Time (TBT)}: Tổng thời gian 
    trình duyệt bị chặn bởi các tác vụ JavaScript dài (trên 50ms), 
    khiến người dùng không thể tương tác.
    
    \item \textbf{Cumulative Layout Shift (CLS)}: Mức độ thay đổi 
    bố cục khi trang tải, ví dụ văn bản bị đẩy xuống do ảnh 
    load chậm.
    
    \item \textbf{Speed Index}: Tốc độ trung bình mà nội dung 
    hiển thị trên màn hình trong quá trình tải trang, chỉ số 
    càng thấp càng tốt.
\end{itemize}

% ─────────────────────────────────────────
\subsection{Kiểm thử chức năng giao diện}

Kiểm thử giao diện được thực hiện trên các chức năng chính của 
hệ thống, mỗi nhóm chức năng chọn một màn hình đại diện.

\begin{figure}[H]
    \centering
    \includegraphics[width=\textwidth]{graphics/test/ui_login.png}
    \caption{Kiểm thử giao diện đăng nhập}
    \label{fig:ui_login}
\end{figure}

\begin{figure}[H]
    \centering
    \includegraphics[width=\textwidth]{graphics/test/ui_product.png}
    \caption{Kiểm thử giao diện danh sách sản phẩm}
    \label{fig:ui_product}
\end{figure}

\begin{figure}[H]
    \centering
    \includegraphics[width=\textwidth]{graphics/test/ui_cart.png}
    \caption{Kiểm thử giao diện giỏ hàng và thanh toán}
    \label{fig:ui_cart}
\end{figure}

\begin{figure}[H]
    \centering
    \includegraphics[width=\textwidth]{graphics/test/ui_tryon.png}
    \caption{Kiểm thử giao diện AI Virtual Try-On}
    \label{fig:ui_tryon}
\end{figure}

\begin{figure}[H]
    \centering
    \includegraphics[width=\textwidth]{graphics/test/ui_ar.png}
    \caption{Kiểm thử giao diện AR Try-On}
    \label{fig:ui_ar}
\end{figure}

\begin{figure}[H]
    \centering
    \includegraphics[width=\textwidth]{graphics/test/ui_chatbot.png}
    \caption{Kiểm thử giao diện Chatbot tư vấn}
    \label{fig:ui_chatbot}
\end{figure}

\begin{figure}[H]
    \centering
    \includegraphics[width=\textwidth]{graphics/test/ui_admin.png}
    \caption{Kiểm thử giao diện quản trị hệ thống}
    \label{fig:ui_admin}
\end{figure}

% ─────────────────────────────────────────
\subsection{Kiểm thử trên nhiều thiết bị}

Sử dụng extension Responsive Viewer trên Chrome để kiểm tra tính 
tương thích giao diện trên nhiều loại thiết bị:

\begin{itemize}
    \item \textbf{Điện thoại}: iPhone 14, iPhone 14 Plus, 
    iPhone SE, Samsung Galaxy S8, Galaxy S20
    \item \textbf{Máy tính bảng}: iPad Mini, iPad Air, iPad Pro
    \item \textbf{Máy tính}: Màn hình 1920$\times$1080px, 
    MacBook Pro 15 inch
\end{itemize}

\begin{figure}[H]
    \centering
    \includegraphics[width=\textwidth]{graphics/test/responsive_mobile.png}
    \caption{Kiểm thử giao diện trên thiết bị di động}
    \label{fig:responsive_mobile}
\end{figure}

\begin{figure}[H]
    \centering
    \includegraphics[width=\textwidth]{graphics/test/responsive_tablet.png}
    \caption{Kiểm thử giao diện trên máy tính bảng}
    \label{fig:responsive_tablet}
\end{figure}

Riêng chức năng \textbf{AR Try-On} được kiểm thử trực tiếp trên 
thiết bị thật do yêu cầu quyền truy cập camera thực tế, không 
thể giả lập qua Responsive Viewer.

% ─────────────────────────────────────────
\subsection{Kiểm thử trên nhiều trình duyệt}

Hệ thống được kiểm thử trên các trình duyệt phổ biến nhằm đảm 
bảo tính tương thích và hiển thị nhất quán.

\begin{table}[H]
\centering
\caption{Kết quả kiểm thử trên nhiều trình duyệt}
\label{tab:browser_test}
\begin{tabular}{|c|l|l|c|}
\hline
\textbf{STT} & \textbf{Trình duyệt} & \textbf{Chức năng kiểm tra} 
& \textbf{Kết quả} \\
\hline
1 & Google Chrome & Toàn bộ chức năng & Đạt \\
\hline
2 & Microsoft Edge & Giao diện, mua sắm, chatbot & Đạt \\
\hline
3 & Mozilla Firefox & Giao diện, mua sắm, chatbot & Đạt \\
\hline
\end{tabular}
\end{table}

\textit{Lưu ý: Chức năng AR Try-On sử dụng Snap Camera Kit hoạt 
động tốt nhất trên Google Chrome do yêu cầu hỗ trợ WebGL và 
quyền truy cập camera của trình duyệt.}